% =============================================================================
% CompetitiveExam/CS401/12-pipelining-risc-cisc-answers.tex
% Independent answer key for 12-pipelining-risc-cisc-questions.tex.
%
% All eight questions are original GATE-pattern questions. Answers were
% independently recomputed from the stated assumptions; no answer is copied
% from an unverified PYQ key.
%
% Compile from CompetitiveExam/CS401 on Windows/MiKTeX:
%   $env:TEXINPUTS = "../../Preambles;"
%   $env:BIBINPUTS = "../..;"
%   xelatex -interaction=nonstopmode 12-pipelining-risc-cisc-answers.tex
% =============================================================================

\documentclass[11pt]{article}
\usepackage[margin=1in]{geometry}
% =============================================================================
% Preambles/header.tex — shared LaTeX/Beamer preamble for this project
%
% Use from a Beamer lecture by prepending:
%     \input{header}
% and compiling with TEXINPUTS=../Preambles:$TEXINPUTS (the /compile-latex
% skill does this automatically).
%
% PALETTE CONTRACT. Color names defined here MUST match the names in
% Quarto/theme-template.scss so Beamer and Quarto renderings use the same
% palette. The scripts/check-palette-sync.sh script enforces this.
%
% To customize: edit the HEX values below AND the matching $variable in
% Quarto/theme-template.scss, then run scripts/check-palette-sync.sh.
% =============================================================================

% -----------------------------------------------------------------------------
% Engine detection + encoding
%
% Our /compile-latex skill uses XeLaTeX, where inputenc is unnecessary and
% can produce encoding warnings. Only load inputenc under pdfTeX.
% -----------------------------------------------------------------------------
\usepackage{iftex}
\ifPDFTeX
  \usepackage[utf8]{inputenc}
\fi

\usepackage{amsmath, amssymb, amsthm}
\usepackage{booktabs}
\usepackage{graphicx}

% -----------------------------------------------------------------------------
% PALETTE — mirrors Quarto/theme-template.scss variable names.
% Load xcolor and define colors BEFORE anything that references them
% (notably hyperref, hypersetup, and the Beamer color assignments).
% -----------------------------------------------------------------------------
\usepackage{xcolor}

\definecolor{primary-blue}{HTML}{012169}      % headings, accents
\definecolor{primary-gold}{HTML}{B9975B}      % emphasis, borders
\definecolor{highlight-yellow}{HTML}{F2A900}  % markers, alerts
\definecolor{light-bg}{HTML}{E8EDF5}          % title / section backgrounds
\definecolor{jet}{HTML}{1A1A1A}               % body text

% Semantic colors (used in diagrams and annotations)
\definecolor{positive}{HTML}{15803D}          % good / observed / identified
\definecolor{negative}{HTML}{B91C1C}          % bad / problematic / bias
\definecolor{neutral}{HTML}{525252}           % reference / context

% Highlight accents (match .hi-* classes in the SCSS)
\definecolor{hi-slate}{HTML}{314F4F}
\definecolor{hi-green}{HTML}{2E8B57}
\definecolor{hi-red}{HTML}{C41E3A}

% Semantic aliases so skill templates and snippets can write
% \textcolor{accent}{...} without hard-coding "primary-blue".
\colorlet{accent}{primary-blue}
\colorlet{accent2}{primary-gold}

% -----------------------------------------------------------------------------
% Hyperref — loaded AFTER xcolor + palette so link colors resolve.
% -----------------------------------------------------------------------------
\usepackage{hyperref}
\hypersetup{colorlinks=true, linkcolor=primary-blue, urlcolor=primary-blue,
            citecolor=primary-blue, pdfborder={0 0 0}}

% -----------------------------------------------------------------------------
% Course code — set once per deck (after \input{header}, before \begin{document})
% via \coursecode{CS401}. Shown in the footer on every slide so a deck's course
% is identifiable without opening the title page. Empty by default.
% -----------------------------------------------------------------------------
\makeatletter
\newcommand{\coursecode}[1]{\gdef\@coursecodetext{#1}}
\gdef\@coursecodetext{}
\makeatother

% -----------------------------------------------------------------------------
% Beamer theme assignments (only apply under Beamer).
%
% We detect Beamer via \@ifclassloaded{beamer}, which is the documented,
% non-internal way. This must be wrapped in \makeatletter/\makeatother
% because \@ifclassloaded contains an @-catcode character.
% -----------------------------------------------------------------------------
\makeatletter
\@ifclassloaded{beamer}{%
  \usefonttheme{professionalfonts}
  \setbeamercolor{structure}{fg=primary-blue}
  \setbeamercolor{frametitle}{fg=primary-blue}
  \setbeamercolor{title}{fg=primary-blue}
  \setbeamercolor{subtitle}{fg=primary-gold}
  \setbeamercolor{author}{fg=primary-blue}
  \setbeamercolor{institute}{fg=neutral}
  \setbeamercolor{date}{fg=primary-gold}
  \setbeamercolor{itemize item}{fg=primary-blue}
  \setbeamercolor{itemize subitem}{fg=primary-gold}
  \setbeamercolor{alerted text}{fg=primary-gold}
  \setbeamercolor{block title}{fg=white, bg=primary-blue}
  \setbeamercolor{block body}{bg=light-bg}
  % Semantic block variants: block=definition/concept (blue), exampleblock=
  % worked example (green/positive), alertblock=key takeaway or caution (gold).
  % Keep to <=2 colored boxes per slide across ALL three variants (INV-7).
  \setbeamercolor{block title example}{fg=white, bg=positive}
  \setbeamercolor{block body example}{bg=light-bg}
  \setbeamercolor{block title alerted}{fg=white, bg=primary-gold}
  \setbeamercolor{block body alerted}{bg=light-bg}

  % Minimal footer: course code (left) + page X / N (right), no navigation clutter
  \setbeamertemplate{navigation symbols}{}
  \setbeamertemplate{footline}{%
    \color{neutral}\footnotesize\hspace{0.7em}\@coursecodetext\hfill
    \insertframenumber/\inserttotalframenumber\hspace{0.7em}\vspace{0.3em}}
}{}
\makeatother

% -----------------------------------------------------------------------------
% TikZ setup — libraries the snippet gallery uses
% -----------------------------------------------------------------------------
\usepackage{tikz}
\usetikzlibrary{arrows.meta, positioning, calc, decorations.pathreplacing,
                fit, shapes.geometric, backgrounds}

% Shared TikZ styles — snippets and hand-written diagrams can pick these up
% via `\begin{tikzpicture}[every node/.style={...}]` or by naming specific
% styles (e.g., `\node[dag-node] (X) at (0,0) {X};`).
\tikzset{
  % Node styles for causal DAGs and similar
  dag-node/.style = {
    draw=primary-blue, fill=light-bg, rounded corners=2pt,
    minimum width=1.2cm, minimum height=0.9cm, align=center, font=\small
  },
  decision-node/.style = {
    draw=primary-gold, fill=white, diamond, aspect=1.8,
    minimum width=1.8cm, minimum height=1.2cm, align=center, font=\small,
    inner sep=1pt
  },
  % Edge styles — observed vs counterfactual
  observed-edge/.style = {-{Latex[length=2.5mm]}, thick, draw=jet},
  counterfactual-edge/.style = {-{Latex[length=2.5mm]}, thick, draw=neutral, dashed},
  confound-edge/.style = {-{Latex[length=2.5mm]}, thick, draw=negative, dashed},
  % Data-point styles for plots
  observed-dot/.style = {fill=primary-blue, draw=primary-blue, circle, inner sep=1.2pt},
  counterfactual-dot/.style = {fill=white, draw=neutral, circle, inner sep=1.2pt}
}

% -----------------------------------------------------------------------------
% Convenience macros
% -----------------------------------------------------------------------------
\newcommand{\muted}[1]{\textcolor{neutral}{#1}}
\newcommand{\key}[1]{\textcolor{primary-gold}{\textbf{#1}}}
\newcommand{\good}[1]{\textcolor{positive}{#1}}
\newcommand{\bad}[1]{\textcolor{negative}{#1}}

% Standout frame macro: full-bleed dark background for section transitions.
% Beamer-only (uses \begin{frame}/\setbeamercolor/\usebeamerfont) -- do not
% call from an article-class file (e.g. Notes/<CODE>/*-notes.tex). \muted,
% \key, \good, \bad, and \coursecode above are plain xcolor/gdef and are
% safe to use from either class; verified when Notes/article-class support
% was added.
% Expand: \transitionslide{Your transition title}
\newcommand{\transitionslide}[1]{%
  \begingroup
  \setbeamercolor{background canvas}{bg=primary-blue}
  \begin{frame}[plain]
    \centering
    \vfill
    {\usebeamerfont{title}\color{white}\Large #1\par}
    \vfill
  \end{frame}
  \endgroup
}

% Section divider frame (Beamer-only), an alternative to \transitionslide with a
% darker two-line style: near-black (jet) background, a small white label line
% above a larger gold title, and a thin gold rule along the bottom edge.
% Expand: \sectiondivider{Section 2}{Pointers}
\newcommand{\sectiondivider}[2]{%
  \begingroup
  \setbeamercolor{background canvas}{bg=jet}
  \begin{frame}[plain]
    \vspace*{3.0cm}
    {\color{white}\ttfamily\large #1\par}
    \vspace{0.5cm}
    {\color{primary-gold}\LARGE #2\par}
    \begin{tikzpicture}[remember picture, overlay]
      \draw[primary-gold, line width=2.5pt]
        ($(current page.south west)+(0,0.1)$) -- ($(current page.south east)+(0,0.1)$);
    \end{tikzpicture}
  \end{frame}
  \endgroup
}

% -----------------------------------------------------------------------------
% Channel watermark — "Concepts That Click" (YouTube).
%
% Article-class files (CompetitiveExam Q/A sets, Notes, Assignments) get a
% light diagonal draft watermark on every page plus a centered footer naming
% the channel. Beamer decks get a subtle TikZ background watermark instead
% (the footline already carries the course code). Centralized here so every
% MCQ PDF is branded without per-file edits.
% -----------------------------------------------------------------------------
\makeatletter
\@ifclassloaded{article}{%
  \usepackage{draftwatermark}
  \SetWatermarkText{Concepts That Click}
  \SetWatermarkScale{1}
  \SetWatermarkFontSize{2cm}
  \SetWatermarkLightness{0.86}
  \SetWatermarkAngle{45}
  \usepackage{fancyhdr}
  \pagestyle{fancy}
  \fancyhf{}
  \fancyfoot[C]{\footnotesize\textcolor{neutral}{Concepts That Click~|~YouTube: \href{https://www.youtube.com/@concepts.thatclick}{@concepts.thatclick}}}
  \renewcommand{\headrulewidth}{0pt}
  \renewcommand{\footrulewidth}{0pt}
  \fancypagestyle{plain}{%
    \fancyhf{}%
    \fancyfoot[C]{\footnotesize\textcolor{neutral}{Concepts That Click~|~YouTube: \href{https://www.youtube.com/@concepts.thatclick}{@concepts.thatclick}}}%
    \renewcommand{\headrulewidth}{0pt}%
    \renewcommand{\footrulewidth}{0pt}%
  }
}{}
\@ifclassloaded{beamer}{%
  \setbeamertemplate{background}{%
    \begin{tikzpicture}[remember picture, overlay]
      \node[rotate=45, opacity=0.055, align=center]
        at (current page.center)
        {\Huge\textcolor{neutral}{Concepts That Click}};
    \end{tikzpicture}%
  }
}{}
\makeatother 
\coursecode{CS401}
\renewcommand{\thesection}{12.\arabic{section}}

\title{Competitive Exam Practice 12 --- Answer Key\\
  \large Instruction Pipelining and RISC vs. CISC}
\author{Auqib Hamid Lone\\[0.3em]
  \emph{Department of Computer Science \& Engineering}, \emph{GCET Ganderbal}}
\date{\today}

\begin{document}
\maketitle

\begin{sloppypar}
\noindent All eight items are \textit{[Original, GATE-pattern]}; the local
scanned PYQ source did not provide a complete, independently verifiable stem
and answer for a Week 12 item. The calculations below are independent checks
of the question file, especially the pipeline schedules and performance
equation.
\end{sloppypar}

\begin{enumerate}
\item[\textbf{Q1}] \textit{[Original, GATE-pattern]} \textbf{10 cycles,
speedup $3.75$.} For $k=5$ and $n=12$,
\[
  T_{\mathrm{pipe}}=k+n-1=5+11=16,
  \qquad T_{\mathrm{seq}}=nk=12\times5=60.
\]
Thus $S=60/16=3.75$. The limiting large-stream speedup is 5, but it is not
the answer for this finite stream.

\item[\textbf{Q2}] \textit{[Original, GATE-pattern]} \textbf{cycle 8,
1 bubble.} The schedule is:
\begin{center}
\begin{tabular}{c|cccccccc}
\textbf{Instruction} & 1 & 2 & 3 & 4 & 5 & 6 & 7 & 8 \\
\hline
\texttt{LW}  & IF & ID & EX & MEM & WB & & & \\
\texttt{ADD} & & IF & ID & stall & EX & MEM & WB & \\
\texttt{SUB} & & & IF & stall & ID & EX & MEM & WB
\end{tabular}
\end{center}
The load-use RAW dependence needs one bubble because the load value is not
available for the immediately following EX stage. The ADD ALU result is
forwarded to SUB, so no second bubble is needed.

\item[\textbf{Q3}] \textit{[Original, GATE-pattern]} \textbf{(B).} A single
instruction still traverses IF, ID, EX, MEM, and WB, so its ideal latency is
about five cycles. Once the pipeline is full, independent instructions can
complete at one per cycle.

\item[\textbf{Q4}] \textit{[Original, GATE-pattern]} \textbf{0.06 CPI.}
The expected penalty per instruction is
\[
  \Delta\mathrm{CPI}=f_{br}p_{miss}P_{flush}
  =0.20\times0.10\times3=0.06.
\]

\item[\textbf{Q5}] \textit{[Original, GATE-pattern]} \textbf{(B).} Two
stages demand one unavailable resource in the same cycle, which is a
structural hazard. Separate instruction and data memories remove this
particular conflict; stalling is another possible remedy.

\item[\textbf{Q6}] \textit{[Original, GATE-pattern]} \textbf{(A), (B), and
(D).} These are design tendencies and a valid measurement rule. Statement (C)
is too strong: instruction count, CPI, clock rate, stalls, cache behavior,
code size, and the workload interact through
$T_{CPU}=\mathrm{IC}\times\mathrm{CPI}/f$.

\item[\textbf{Q7}] \textit{[Original, GATE-pattern]} \textbf{(B).} Using
$T_{CPU}=\mathrm{IC}\times\mathrm{CPI}/f$,
\[
  T_R=\frac{1.20\times10^9\times1.1}{2.4\times10^9}=0.550~\mathrm{s},
  \qquad
  T_C=\frac{0.80\times10^9\times2.0}{3.0\times10^9}
      =0.533\overline{3}~\mathrm{s}.
\]
The CISC-like implementation is faster for this stated workload despite its
higher CPI.

\item[\textbf{Q8}] \textit{[Original, GATE-pattern]} \textbf{(A), (B), and
(C).} A single instruction over many data elements is SIMD-style; independent
cores running independent instruction streams are MIMD-style. A large working
set can expose cache and virtual-memory costs. Statement (D) is false because
total time also depends on instruction count, CPI, stalls, and memory behavior.
\end{enumerate}

\section*{Source Boundaries and Verification}
The questions are derived from the completed Week 12 deck and notes. The
textbook concepts are standard treatments supported by the indexed sources:
Patterson and Hennessy, ARM 5e, Ch.~4 §§4.5--4.8 (pipeline and hazards),
Ch.~6 §§6.2--6.3 (parallelism), and App.~D (RISC survey); Stallings 11e,
Ch.~16 §16.4 and Ch.~17 §§17.5, 17.9 (pipeline and RISC/CISC). These are
supporting references, not claims that any constructed question appeared in
those books or in GATE. The local PYQ registry identifies Week 12 pages, but
OCR did not recover a complete stem and answer; therefore no real PYQ claim is
made here. All eight answers were recomputed from the explicit assumptions in
the stems, including the two timing tables and both CPI/time calculations.

\end{document}